% =====================================================================
%  مذكّرة تخرّج (ليسانس) — دراسة نوعية
%  إدراك أسلوب المعاملة الوالدية المتسلطة وعلاقته بتقدير الذات لدى المراهقين
%
%  المحرّك المطلوب: XeLaTeX  (وليس pdfLaTeX)
%  يُترجم مباشرةً على Overleaf:  Menu > Compiler > XeLaTeX
%  الخطّ المستعمل: Amiri (متوفّر على Overleaf)
% =====================================================================
\documentclass[12pt,oneside]{report}

% =====================================================================
%  ترتيب الحِزَم مهمّ: يجب تحميل bidi (الذي تجلبه polyglossia للعربية)
%  بعد hyperref و pgfplots. لذلك نُحمّل اللغة العربية في آخر الديباجة.
% =====================================================================

% ------------------------- الخطوط (fontspec أولاً) -------------------
\usepackage{fontspec}

% ------------------------- حِزَم عامّة --------------------------------
\usepackage{geometry}
\geometry{a4paper, top=2.5cm, bottom=2.5cm, left=2.5cm, right=2.5cm}
\usepackage{setspace}
\onehalfspacing
\usepackage{xcolor}
\definecolor{titleblue}{RGB}{31,73,125}
\definecolor{lightgray}{RGB}{240,240,240}
\usepackage{graphicx}
\usepackage{array}
\usepackage{longtable}
\usepackage{booktabs}
\usepackage{colortbl}
\usepackage{multirow}
\usepackage{amssymb}      % للرمز \checkmark
\usepackage{enumitem}

% ------------------------- المخطّطات (charts) ------------------------
\usepackage{tikz}
\usepackage{pgfplots}
\pgfplotsset{compat=1.18}
\usetikzlibrary{patterns}

% ------------------------- الروابط -----------------------------------
\usepackage{hyperref}
\hypersetup{colorlinks=true, linkcolor=titleblue, urlcolor=titleblue, unicode=true}

% ------------------------- عناوين الفصول -----------------------------
\usepackage{titlesec}
\titleformat{\chapter}[display]
  {\normalfont\huge\bfseries\color{titleblue}}
  {\chaptertitlename\ \thechapter}{14pt}{\Huge}
\titlespacing*{\chapter}{0pt}{10pt}{30pt}

% تسمية الأقسام بالعربية
\renewcommand{\contentsname}{فهرس المحتويات}
\renewcommand{\listtablename}{فهرس الجداول}
\renewcommand{\listfigurename}{فهرس الأشكال}
\renewcommand{\chaptername}{الفصل}
\renewcommand{\tablename}{جدول}
\renewcommand{\figurename}{شكل}
\renewcommand{\bibname}{قائمة المراجع}

% أمر مختصر للمراجع اللاتينية داخل النصّ العربي
\newcommand{\en}[1]{\textenglish{#1}}
% عنصر مرجع بنمط APA مع إزاحة معلّقة
\newcommand{\refitem}[1]{\par\noindent\hangindent=1.8em\hangafter=1 #1\par\vspace{4pt}}

\setcounter{tocdepth}{2}
\setcounter{secnumdepth}{2}

% ------------------------- اللغة العربية (آخر شيء) -------------------
\usepackage{polyglossia}            % تجلب bidi وتُحمَّل بعد hyperref/pgfplots
\setmainlanguage[numerals=maghrib]{arabic} % أرقام أوروبية (1234567890)
\setotherlanguage{english}
\newfontfamily\arabicfont[Script=Arabic]{Amiri}
\newfontfamily\englishfont{Latin Modern Roman}
% بيئة مختصرة لرسم المخطّطات في اتّجاه LTR لتفادي انعكاس المحاور
\newcommand{\chartLTR}[1]{\begin{LTR}#1\end{LTR}}

% =====================================================================
\begin{document}

% =====================================================================
%                            صفحة الغلاف
% =====================================================================
\begin{titlepage}
\centering
{\large الجمهورية الجزائرية الديمقراطية الشعبية}\\[2pt]
{\large وزارة التعليم العالي والبحث العلمي}\\[14pt]
{\large جامعة ....................}\\[2pt]
{\large كلية العلوم الإنسانية والاجتماعية}\\[2pt]
{\large قسم علم النفس}\\[40pt]

\rule{\linewidth}{0.5pt}\\[10pt]
{\Huge\bfseries\color{titleblue} إدراك أسلوب المعاملة الوالدية المتسلطة\\[10pt]
وعلاقته بتقدير الذات لدى المراهقين}\\[12pt]
{\Large دراسة نوعية لعيّنة من المراهقين (15--16 سنة)}\\[8pt]
\rule{\linewidth}{0.5pt}\\[30pt]

{\large مذكّرة مكمّلة لنيل شهادة الليسانس في علم النفس}\\[40pt]

\begin{minipage}{0.45\linewidth}
\centering
{\bfseries من إعداد الطالب(ة):}\\[6pt]
.............................
\end{minipage}\hfill
\begin{minipage}{0.45\linewidth}
\centering
{\bfseries إشراف الأستاذ(ة):}\\[6pt]
.............................
\end{minipage}\\[60pt]

{\large السنة الجامعية: 2025 / 2026}
\end{titlepage}

% =====================================================================
%                       شكر  +  ملخّص  +  الفهارس
% =====================================================================
\chapter*{شكر وتقدير}
\addcontentsline{toc}{chapter}{شكر وتقدير}
الحمد لله الذي وفّقنا لإتمام هذا العمل المتواضع. نتقدّم بخالص الشكر والتقدير إلى الأستاذ(ة)
المشرف(ة) على ما قدّمه من توجيهات قيّمة طوال مراحل إنجاز هذه المذكّرة، وإلى كلّ من ساهم في
إتمام هذا البحث، وفي مقدّمتهم المراهقون المشاركون الذين منحونا ثقتهم ووقتهم.

\chapter*{ملخّص الدراسة}
\addcontentsline{toc}{chapter}{ملخّص الدراسة}
هدفت هذه الدراسة النوعية إلى الكشف عن الكيفية التي يُدرك بها المراهقون (15--16 سنة) أسلوب
المعاملة الوالدية المتسلطة، وانعكاس هذا الإدراك على تقديرهم لذاتهم. اعتمدت الدراسة على المنهج
الوصفي بأسلوبه الكيفي، واستُخدمت المقابلة النصف موجَّهة أداةً لجمع البيانات، طُبِّقت على عيّنة
قصدية مكوّنة من تسعة (09) مراهقين. وقد عُولجت البيانات عن طريق تحليل المضمون. وأظهرت النتائج
أنّ المراهقين يُدركون الأسلوب الوالدي المتسلط من خلال غياب النقاش (100\%) والعقاب اللفظي
(100\%) وغياب الدفء العاطفي (89\%)، وأنّ هذا الإدراك يقترن بتقدير ذات منخفض تجلّى في الشعور
بعدم الكفاءة (89\%) والقيمة المشروطة المنخفضة (89\%)، بمتوسّط درجة تقدير ذاتٍ بلغ 4.11 من 10.
كما كشفت الدراسة عن حالات أبدت تكيّفاً ظاهرياً (33\%) ارتبط بإدراكٍ أقلّ حدّةً للتسلّط.\\[6pt]
\noindent\textbf{الكلمات المفتاحية:} إدراك، أسلوب المعاملة الوالدية المتسلطة، تقدير الذات،
المراهقة، الدراسة النوعية.

\tableofcontents
\listoftables
\listoffigures

% =====================================================================
%                              المقدّمة
% =====================================================================
\chapter*{مقدّمة}
\addcontentsline{toc}{chapter}{مقدّمة}

تُعدّ الأسرة الخليّة الأولى التي ينشأ فيها الفرد ويتشكّل من خلالها جانب كبير من بنيته النفسية
والاجتماعية، ويحتلّ أسلوب المعاملة الوالدية مكانة محورية في توجيه مسار النمو النفسي للأبناء،
لا سيّما في مرحلة المراهقة التي تُمثّل منعطفاً حاسماً في بناء الهوية وتقدير الذات. ومن بين
الأساليب الوالدية التي حظيت باهتمام واسع في الأدبيات النفسية، يبرز الأسلوب المتسلط بوصفه نمطاً
يقوم على الضبط الصارم والطاعة المطلقة وغياب الحوار، وهو ما قد ينعكس سلباً على الصورة التي
يُكوّنها المراهق عن ذاته.

وقد جاءت هذه الدراسة لتستكشف، بأسلوب كيفي، الكيفية التي \emph{يُدرك} بها المراهق هذا الأسلوب
الوالدي، وكيف يتمثّله ويُفسّره، وما ينعكس عنه على تقديره لذاته، انطلاقاً من تجربته الذاتية لا
من قياس كمّي خارجي.

وقد قُسّمت هذه المذكّرة، تبعاً للمتعارف عليه، إلى أربعة فصول يسبقها هذا التقديم؛ خُصّص
\textbf{الفصل الأول} للإطار النظري للمفهوم الأول (إدراك أسلوب المعاملة الوالدية المتسلطة) من
حيث تقديمه وتعريفاته وأبعاده وتعريفه الإجرائي. أمّا \textbf{الفصل الثاني} فعالج المفهوم الثاني
(تقدير الذات لدى المراهقين) بالطريقة نفسها. وتناول \textbf{الفصل الثالث} الإطار المنهجي للدراسة
من خلال بسط الإشكالية والتساؤلات والدراسات السابقة والفجوة البحثية وأهمية البحث وأهدافه، ثمّ
وصف المنهج والعيّنة والأداة وشبكة المقابلة. وأُفرد \textbf{الفصل الرابع} لعرض النتائج وتحليلها
إحصائياً (وصفياً وتحليلياً) ثمّ مناقشتها في ضوء الدراسات السابقة، لتُختتم المذكّرة بخلاصة عامّة
وقائمة المراجع والملاحق.

% =====================================================================
%                  الفصل الأول — الإطار النظري (المفهوم 1)
% =====================================================================
\chapter{الإطار النظري: إدراك أسلوب المعاملة الوالدية المتسلطة}

\section{تمهيد}
يُعدّ أسلوب المعاملة الوالدية من أكثر المتغيّرات النفسية والاجتماعية تأثيراً في تشكيل شخصية
الطفل وتطوّر مفهوم الذات لديه خلال مرحلة المراهقة. وتتعدّد هذه الأساليب وتتباين تبعاً لمنظومة
القيم الثقافية والاجتماعية السائدة في كلّ مجتمع. غير أنّ الأدبيات النفسية تتّفق على أنّ الأسلوب
الوالدي المتسلط يُمثّل نمطاً من أكثر الأنماط إثارةً للجدل لما يُفرزه من تداعيات على الأبناء
المراهقين. ويُقصد بإدراك أسلوب المعاملة الوالدية المتسلطة الطريقة التي يُفسّر بها المراهق
ويعي ويُقيّم السلوكيات الوالدية الضابطة التي يتلقّاها في محيطه الأسري.

\section{التعريفات النظرية لأسلوب المعاملة الوالدية المتسلطة}

\subsection{تعريف ديانا بومرين (Baumrind)}
تُعدّ ديانا بومرين الرائدة الأولى في تصنيف أساليب التنشئة الوالدية، إذ ميّزت بين ثلاثة أساليب
أساسية هي: المتسلط \en{(Authoritarian)}، والسلطوي التشاركي \en{(Authoritative)}، والمتساهل
\en{(Permissive)}. وتُعرّف بومرين الأسلوب المتسلط بأنّه نمط من التحكّم يسعى إلى تشكيل سلوك
الطفل وضبطه وتقييمه وفق معايير ثابتة وصارمة يعتبرها الوالدان مطلقة وغير قابلة للتفاوض، مع
التأكيد على الطاعة وتقييد الاستقلالية واستخدام العقوبة، وغياب شبه تامّ للتفسير والحوار
\en{(Baumrind, 1966)}. وأضافت لاحقاً أنّ هذا الأسلوب يتّسم بمستوى عالٍ من الضبط والمطالب
السلوكية يقابله مستوى منخفض من الاستجابة الوجدانية والدفء العاطفي \en{(Baumrind, 1991)}.

\subsection{تعريف ماكوبي ومارتن (Maccoby \& Martin)}
طوّر ماكوبي ومارتن نموذج بومرين بإضافة بُعد رابع هو الأسلوب المُهمِل، وعرّفا الأسلوب المتسلط
بوصفه نمطاً يتّسم بمستوى مرتفع من المطالب والتحكّم ومستوى منخفض من الاستجابة الوجدانية، مؤكّدَين
أنّ الضبط الصارم في غياب الدفء العاطفي يُفضي إلى توتّرات نفسية لدى الأبناء، لا سيّما في مرحلة
المراهقة \en{(Maccoby \& Martin, 1983)}.

\subsection{تعريف باربر (Barber)}
يُميّز باربر بين نوعين من التحكّم الوالدي في سياق الأسلوب المتسلط: \textbf{التحكّم السلوكي}
\en{(Behavioral Control)} ويعني مراقبة سلوكيات الأبناء الخارجية، و\textbf{التحكّم النفسي}
\en{(Psychological Control)} ويشير إلى التحكّم في أفكار المراهق ومشاعره وهويته الداخلية، وهو
الأكثر ضرراً على الصحّة النفسية \en{(Barber, 1996)}.

\section{أبعاد أسلوب المعاملة الوالدية المتسلطة}
استناداً إلى نموذج بومرين والأطر المكمّلة، يمكن تحديد الأبعاد الرئيسية في:
\begin{description}[leftmargin=1.5em, style=nextline]
  \item[البُعد الأول — الضبط والسيطرة الصارمة:] فرض قواعد صارمة غير مرنة وتطبيقها دون نقاش، فيُدرك
  المراهق غياب حرية الاختيار وثقل القيود \en{(Baumrind, 1966)}.
  \item[البُعد الثاني — غياب الدفء العاطفي:] ندرة التعبير الوالدي عن الحبّ والدعم، فيُدرك المراهق
  العلاقة بوصفها علاقة سلطة لا علاقة محبّة واحتواء \en{(Maccoby \& Martin, 1983)}.
  \item[البُعد الثالث — استخدام العقوبة:] اللجوء إلى أساليب تأديبية قسرية جسدية أو نفسية كالإهانة
  والحرمان والتهديد، فيقترن إدراك المراهق لها بالخوف والقلق.
  \item[البُعد الرابع — تقييد الاستقلالية:] منع المراهق من اتّخاذ قراراته والتعبير عن رأيه، ممّا
  يُعيق تطوّر الهوية المستقلّة \en{(Barber, 1996)}.
\end{description}

\section{التعريف الإجرائي}
يُعرَّف إدراك أسلوب المعاملة الوالدية المتسلطة إجرائياً في هذه الدراسة بأنّه: \emph{مجموع
التمثّلات المعرفية والانفعالية التي يبنيها المراهق (15--16 سنة) حول طريقة تعامل والديه معه،
والتي تتجلّى في إدراكه لأربعة أبعاد هي: الضبط الصارم، وغياب الدفء العاطفي، واستخدام العقوبة،
وتقييد الاستقلالية}، ويُقاس ذلك من خلال استجاباته على بنود شبكة المقابلة النصف موجَّهة المُعدّة
لأغراض هذه الدراسة.

% =====================================================================
%                  الفصل الثاني — الإطار النظري (المفهوم 2)
% =====================================================================
\chapter{الإطار النظري: تقدير الذات لدى المراهقين}

\section{تمهيد}
يُمثّل تقدير الذات أحد المفاهيم المحورية في علم النفس، إذ يرتبط ارتباطاً وثيقاً بالصحّة النفسية
للفرد وتوافقه الشخصي والاجتماعي، ويتّخذ أهمية بالغة خلال مرحلة المراهقة التي يخوض فيها الفرد
تحوّلات جذرية ويسعى إلى بناء هويته وتحقيق ذاته. ويتشكّل تقدير الذات لدى المراهق من خلال تفاعل
عوامل متعدّدة يأتي في مقدّمتها نوعية العلاقات الأسرية وطبيعة المعاملة الوالدية.

\section{التعريفات النظرية لتقدير الذات}

\subsection{تعريف ماسلو (Maslow)}
أدرج ماسلو تقدير الذات ضمن هرمية الحاجات الإنسانية في مستواها الرابع، مُميّزاً بين تقدير الذات
الداخلي (الكفاءة والإنجاز والاستقلالية) وتقدير الذات الخارجي (الاحترام والاعتراف من الآخرين)،
مؤكّداً أنّ إشباع هذه الحاجة شرط للوصول إلى تحقيق الذات \en{(Maslow, 1954)}.

\subsection{تعريف روزنبرغ (Rosenberg)}
يُعرّف روزنبرغ تقدير الذات بأنّه الاتّجاه الإيجابي أو السلبي نحو الذات بوصفها موضوعاً كلّياً،
مؤكّداً أنّ تقدير الذات المرتفع لا يعني الغرور بل يعني أن يعتبر الفرد نفسه ذا قيمة وإن اعترف
بنقائصه، في حين يعني تقدير الذات المنخفض رفض الفرد لذاته وعدم رضاه عنها \en{(Rosenberg, 1965)}.
وقد طوّر مقياسه الشهير الذي لا يزال من أكثر أدوات القياس استخداماً عالمياً.

\subsection{تعريف كوبر سميث (Coopersmith)}
يُعرّف كوبر سميث تقدير الذات بأنّه التقييم المعتاد والمستمرّ الذي يُجريه الفرد لذاته ويُعبّر عن
اتّجاه الموافقة أو الرفض تجاه نفسه ومدى اعتقاده بأنّه كفء وذو أهمية وناجح وقيّم، مُميّزاً بين
أربعة أبعاد: الكفاءة الاجتماعية والأكاديمية والأسرية والصورة العامّة للذات \en{(Coopersmith, 1967)}.

\subsection{تعريف هارتر (Harter)}
طوّرت هارتر نموذجاً خاصاً لتقدير الذات في المراهقة مُميّزةً بين عدّة مجالات: المظهر الجسمي
والقبول الاجتماعي والكفاءة الأكاديمية والسلوك الأخلاقي والتقدير العامّ، مؤكّدةً أنّ الفجوة بين
ما يطمح إليه المراهق وما يُدركه واقعاً هي التي تُحدّد مستوى تقديره العامّ لذاته \en{(Harter, 1988)}.

\section{أبعاد تقدير الذات لدى المراهقين}
\begin{description}[leftmargin=1.5em, style=nextline]
  \item[الكفاءة الذاتية المُدرَكة:] إيمان المراهق بقدرته على مواجهة التحدّيات وإنجاز المهام، وهو
  وثيق الصلة بمفهوم الفاعلية الذاتية لدى \en{Bandura (1977)}.
  \item[الصورة الجسمية:] تكتسب أهمية قصوى في المراهقة المبكرة (15--16 سنة) لما يصاحبها من تحوّلات
  جسمية سريعة ترتبط مباشرةً بتقدير الذات العامّ.
  \item[القبول الاجتماعي:] شعور المراهق بالانتماء والقبول من الأقران، وهو مؤشّر حسّاس في هذه المرحلة.
  \item[الكفاءة الأكاديمية:] إدراك المراهق لقدرته على النجاح الدراسي وتحقيق التوقّعات الأسرية والمدرسية.
  \item[التقدير العامّ للذات:] الشعور الكلّي بالقيمة والجدارة بصرف النظر عن المجالات الجزئية.
\end{description}

\section{خصائص تقدير الذات في مرحلة المراهقة (15--16 سنة)}
تُعدّ هذه المرحلة فارقة في تطوّر تقدير الذات؛ إذ يعيش فيها المراهق إعادة بناء الهوية وفق ما
وصفه إريكسون في مفهوم «أزمة الهوية مقابل التشتّت» \en{(Erikson, 1968)}، فيكون تقدير الذات في
حالة سيولة وتحوّل. كما يُصبح المراهق أكثر وعياً بالتناقض بين «الذات الحقيقية» و«الذات المثالية»،
وأكثر حساسيةً للنقد والرقابة الخارجية، في وقتٍ يبلغ فيه تفاعل الضغط الوالدي والضغط الاجتماعي
ذروته، ممّا يجعل البيئة الأسرية ذات أثر بالغ في رسم ملامح تقديره لذاته.

\section{التعريف الإجرائي}
يُعرَّف تقدير الذات إجرائياً في هذه الدراسة بأنّه: \emph{التقييم الذاتي الذي يحمله المراهق
(15--16 سنة) تجاه قيمته وكفاءته وجدارته كشخص، ويتجلّى في إدراكه لذاته عبر خمسة أبعاد: الكفاءة
الذاتية المُدرَكة، والصورة الجسمية، والقبول الاجتماعي، والكفاءة الأكاديمية، والتقدير العامّ
للذات}، ويُقاس من خلال تحليل مضمون استجاباته في المقابلة النصف موجَّهة.

% =====================================================================
%                     الفصل الثالث — الإطار المنهجي
% =====================================================================
\chapter{الإطار المنهجي للدراسة}

\section{بسط الإشكالية}
تُشكّل الأسرة البيئة الأولى والأكثر تأثيراً في تشكيل الشخصية، ويحتلّ الأسلوب الوالدي المتسلط
مكانةً بارزة في السياق العربي والمغاربي لتجذّره في الموروث الثقافي والاجتماعي. وعلى الرغم من
التراكم المعرفي الكبير في الأدبيات الغربية حول علاقة هذا الأسلوب بتقدير الذات، فإنّ السياق
المغاربي لا يزال يفتقر إلى دراسات \emph{نوعية} تستكشف الكيفية التي يُدرك بها المراهق العربي هذا
الأسلوب وما ينعكس عنه على تقديره لذاته. ولمّا كانت المراهقة (15--16 سنة) مرحلة بالغة الحساسية
يسعى فيها الفرد إلى تأكيد ذاته في مواجهة السلطة الوالدية، فإنّ تقدير الذات فيها يكون هشّاً
ومعرّضاً للتأثّر المزدوج بالبيئة الأسرية والمحيط الاجتماعي.

\section{الدراسات السابقة}
\subsection{الدراسة الأجنبية}
في دراسة أُجريت بالولايات المتحدة على عيّنة تجاوزت أربعة آلاف مراهق، توصّل \en{Lamborn et al.
(1991)} إلى أنّ المراهقين المنتمين إلى أسرٍ ذات أسلوب متسلط يُسجّلون مستويات منخفضة من تقدير
الذات مقارنةً بأقرانهم من أسرٍ ذات أسلوب سلطوي تشاركي، وأنّ الارتباط بين التحكّم الصارم وانخفاض
تقدير الذات يتعزّز خلال المراهقة المبكرة.

\subsection{الدراسة العربية}
وفي الجزائر، توصّل \en{بن عيسى وبلعيد (2018)} إلى وجود علاقة عكسية ذات دلالة إحصائية بين الأسلوب
الوالدي المتسلط وتقدير الذات لدى المراهقين الجزائريين، مع إشارتهما إلى أنّ الخصوصية الثقافية
المغاربية قد تُعدّل من حدّة هذه العلاقة في بعض الحالات.

\subsection{الفجوة البحثية}
رغم وفرة الدراسات الكمّية التي أثبتت العلاقة بين المتغيّرَين، يظلّ هناك شُحّ في الدراسات النوعية
التي تستكشف الكيفية الذاتية التي \emph{يُدرك} بها المراهق المغاربي هذا الأسلوب وأثره على تقدير
ذاته من منظوره الشخصي، وهو ما تسعى هذه الدراسة إلى سدّه.

\section{تساؤلات الدراسة}
\subsection*{السؤال الإشكالي الرئيسي}
\begin{quote}
\textbf{كيف يُدرك المراهقون (15--16 سنة) أسلوب المعاملة الوالدية المتسلطة، وما انعكاس هذا
الإدراك على تقديرهم لذاتهم؟}
\end{quote}
\subsection*{الأسئلة الفرعية}
\begin{enumerate}[leftmargin=2em]
  \item كيف يصف المراهق الضبط الوالدي الصارم الذي يتلقّاه داخل الأسرة؟
  \item كيف يُدرك المراهق غياب الدفء العاطفي في علاقته بوالديه؟
  \item ما الأثر الذي يُدركه المراهق لأسلوب الوالدين المتسلط على كفاءته الذاتية وصورته عن نفسه؟
  \item كيف تنعكس إدراكات المراهق للأسلوب المتسلط على مستوى تقديره لذاته في أبعاده المختلفة؟
\end{enumerate}

\section{أهمية البحث وأهدافه}
\textbf{الأهمية النظرية:} إثراء الأدبيات العربية المتعلّقة بأساليب التنشئة وتقدير الذات في
السياق المغاربي، وتعميق فهم نموذج بومرين في ضوء الخصوصية الثقافية المحلّية.\\
\textbf{الأهمية التطبيقية:} تقديم فهمٍ معمّق يُمكّن المرشدين النفسيين والتربويين من تطوير برامج
تدخّلية تستهدف تعزيز تقدير الذات لدى المراهقين المُتلقّين لأساليب تنشئة متسلطة.\\
\textbf{أهداف البحث:} الكشف عن كيفية إدراك المراهق للأسلوب المتسلط في أبعاده الأربعة، وتحديد
انعكاس هذا الإدراك على أبعاد تقدير الذات، ورصد الحالات المغايرة (التكيّف الظاهري).

\section{منهج الدراسة}
اعتمدت الدراسة على \textbf{المنهج الوصفي بأسلوبه الكيفي (النوعي)} لكونه الأنسب لاستكشاف
التمثّلات الذاتية والمعاني التي يُضفيها المراهق على تجربته مع والديه، بعيداً عن القياس الكمّي.

\section{وصف العيّنة}
\subsection{نوع العيّنة وحجمها}
اعتمدت الدراسة على \textbf{عيّنة قصدية} مكوّنة من \textbf{تسعة (09) مراهقين} تتراوح أعمارهم بين
15 و16 سنة، اختيروا وفق معايير تتلاءم مع طبيعة الدراسة وأهدافها.

\subsection{مبرّرات اختيار الفئة العمرية (15--16 سنة)}
\begin{itemize}[leftmargin=2em]
  \item تُمثّل المراهقة المبكّرة إلى الوسطى المرحلة الأعلى حساسيةً للأسلوب الوالدي.
  \item يشتدّ فيها التعارض بين حاجة المراهق إلى الاستقلالية وما يفرضه الأسلوب المتسلط من قيود.
  \item هي المرحلة الفارقة في تشكّل تقدير الذات وفق نموذجَي هارتر وروزنبرغ.
  \item يُتيح التطوّر المعرفي (التفكير المجرّد) فيها قدرةً أعمق على التأمّل والتعبير عن الإدراكات.
\end{itemize}

\subsection{معايير الاختيار}
\textbf{معايير الشمول:} أن يكون المشارك في الفئة العمرية 15--16 سنة، مقيماً داخل أسرته، يعيش
نمطاً والدياً يتّسم بالضبط الصارم، موافقاً على المشاركة الطوعية.\\
\textbf{معايير الإقصاء:} وجود اضطرابات نفسية مُشخّصة قد تُشوّش النتائج، أو رفض المشاركة.

\subsection{خصائص العيّنة}
\begin{table}[h!]
\centering
\caption{توزيع أفراد العيّنة حسب الجنس والسنّ والمستوى الدراسي}
\renewcommand{\arraystretch}{1.3}
\begin{tabular}{>{\bfseries}c c c c}
\rowcolor{lightgray}
\textnormal{\textbf{الرقم}} & \textbf{الجنس} & \textbf{السنّ} & \textbf{المستوى الدراسي}\\
\hline
م1 & ذكر  & 15 & الثانية ثانوي \\
م2 & أنثى & 16 & الثانية ثانوي \\
م3 & ذكر  & 16 & الثالثة ثانوي \\
م4 & أنثى & 15 & الثانية ثانوي \\
م5 & ذكر  & 15 & الثانية ثانوي \\
م6 & أنثى & 16 & الثالثة ثانوي \\
م7 & ذكر  & 16 & الثالثة ثانوي \\
م8 & أنثى & 15 & الثانية ثانوي \\
م9 & ذكر  & 16 & الثالثة ثانوي \\
\hline
\end{tabular}
\label{tab:sample}
\end{table}

% ---------------- مخطّط 1: الجنس --------------------
\begin{figure}[h!]
\centering
\chartLTR{%
\begin{tikzpicture}
\begin{axis}[
  ybar, width=11cm, height=6cm,
  bar width=28pt, ymin=0, ymax=6,
  ylabel={\textarabic{عدد المشاركين}},
  symbolic x coords={males,females},
  xtick=data,
  xticklabels={\textarabic{ذكور},\textarabic{إناث}},
  nodes near coords, enlarge x limits=0.5,
  ylabel style={font=\arabicfont},
  tick label style={font=\arabicfont},
]
\addplot[fill=titleblue] coordinates {(males,5) (females,4)};
\end{axis}
\end{tikzpicture}}
\caption{توزيع العيّنة حسب الجنس (ذكور = 5، إناث = 4)}
\label{fig:gender}
\end{figure}

% ---------------- مخطّط 2: السنّ --------------------
\begin{figure}[h!]
\centering
\chartLTR{%
\begin{tikzpicture}
\begin{axis}[
  ybar, width=11cm, height=6cm,
  bar width=28pt, ymin=0, ymax=6,
  ylabel={\textarabic{عدد المشاركين}},
  symbolic x coords={a15,a16},
  xtick=data,
  xticklabels={\textarabic{15 سنة},\textarabic{16 سنة}},
  nodes near coords, enlarge x limits=0.5,
  ylabel style={font=\arabicfont},
  tick label style={font=\arabicfont},
]
\addplot[fill=titleblue!70] coordinates {(a15,4) (a16,5)};
\end{axis}
\end{tikzpicture}}
\caption{توزيع العيّنة حسب السنّ (15 سنة = 4، 16 سنة = 5)}
\label{fig:age}
\end{figure}

\section{وصف الأداة: المقابلة النصف موجَّهة}
\subsection{تعريف الأداة وتبريرها}
اعتمدت الدراسة على \textbf{المقابلة النصف موجَّهة} \en{(L'entretien semi-directif)} لكونها الأنسب
لطبيعتها؛ إذ تتيح للباحث التأطير الموضوعاتي للمقابلة، وتمنح المشارك في الوقت ذاته حرية التعبير
عن تجربته وإدراكاته بعيداً عن قيود المقابلة الموجَّهة.

\subsection{الإطار النظري لبناء الأداة}
بُنيت شبكة المقابلة على محورين نظريين: نموذج بومرين \en{(1966, 1991)} في أبعاد الأسلوب المتسلط،
ونموذجَي روزنبرغ \en{(1965)} وهارتر \en{(1988)} في أبعاد تقدير الذات. وصُمّمت الأسئلة لتستكشف
علاقة الإدراك المتبادل بين المتغيّرَين، مع التركيز على المفهوم المحوري: كيف يُدرك المراهق
الأسلوب المتسلط وانعكاسه على تقديره لذاته. وتجدر الإشارة إلى أنّ ما يُسجَّل في الشبكة هو
\emph{الأسئلة} الموزّعة على المحاور، بينما تُدرَج \emph{الإجابات} الكاملة في الملاحق
(انظر الملحق~\ref{app:transcripts}).

\subsection{شبكة المقابلة النصف موجَّهة}
\paragraph{المحور الأول — إدراك الضبط الصارم والسيطرة الوالدية:}
\begin{enumerate}[leftmargin=2em]
  \item كيف تصف العلاقة بينك وبين والديك في الحياة اليومية؟
  \item هل ثمّة قواعد وقيود في بيتك؟ صِف لي كيف تُوضع ومن يضعها.
  \item ماذا يحدث عندما لا تلتزم بالقواعد المفروضة عليك في المنزل؟
  \item هل يمنحك والداك مجالاً لإبداء رأيك في القرارات التي تمسّك؟
\end{enumerate}
\paragraph{المحور الثاني — إدراك غياب الدفء العاطفي:}
\begin{enumerate}[leftmargin=2em, start=5]
  \item كيف تشعر تجاه والديك عندما تُخطئ أو تفشل في أمر ما؟
  \item هل تشعر أنّ والديك يتفهّمان مشاعرك واحتياجاتك النفسية؟
  \item هل يمكنك أن تتحدّث إلى والديك بحرية عمّا يزعجك أو يُقلقك؟
\end{enumerate}
\paragraph{المحور الثالث — إدراك أساليب العقوبة والتأديب:}
\begin{enumerate}[leftmargin=2em, start=8]
  \item كيف يتعامل والداك معك حين تتصرّف بخلاف ما يريدان؟
  \item هل عانيت من أساليب عقاب تجعلك تشعر بالخوف أو الإهانة؟
  \item كيف يؤثّر أسلوب التأديب الذي تتلقّاه على علاقتك بوالديك؟
\end{enumerate}
\paragraph{المحور الرابع — إدراك تقييد الاستقلالية:}
\begin{enumerate}[leftmargin=2em, start=11]
  \item هل تشعر أنّ لديك حرية في اتّخاذ قراراتك كاختيار أصدقائك وهواياتك؟
  \item كيف تتصرّف عندما ترغب في شيء يخالف إرادة والديك؟
  \item كيف تُلقي القيود التي يفرضها والداك بظلالها على شعورك بنفسك؟
\end{enumerate}
\paragraph{المحور الخامس — تقدير الذات: الكفاءة الذاتية المُدرَكة:}
\begin{enumerate}[leftmargin=2em, start=14]
  \item كيف تُقيّم قدرتك على التعامل مع المشكلات والتحدّيات؟
  \item هل تعتقد أنّ طريقة تعامل والديك تؤثّر على ثقتك بنفسك وقدراتك؟
  \item هل تشعر أنّك قادر على اتّخاذ قراراتك وحدك أم تحتاج دائماً إلى إذن أو توجيه؟
\end{enumerate}
\paragraph{المحور السادس — تقدير الذات: الصورة الذاتية والقبول الاجتماعي:}
\begin{enumerate}[leftmargin=2em, start=17]
  \item كيف ترى نفسك بصفة عامّة؟ وما الكلمات التي تصف بها شخصيتك؟
  \item هل تعتقد أنّ طريقة تعامل والديك تؤثّر على تعاملك مع أصدقائك ومحيطك؟
  \item هل تشعر بالقبول والانتماء في وسطك الاجتماعي؟ وهل لأسلوب والديك علاقة بذلك؟
\end{enumerate}
\paragraph{المحور السابع — تقدير الذات العامّ:}
\begin{enumerate}[leftmargin=2em, start=20]
  \item إذا طُلب منك أن تُقيّم قيمتك وجدارتك كشخص من صفر إلى عشرة، فماذا تعطي نفسك ولماذا؟
  \item ما الذي تتمنّى تغييره في طريقة معاملة والديك لك؟
  \item كلمة أخيرة: كيف تُلخّص أثر معاملة والديك على نظرتك إلى نفسك؟
\end{enumerate}

% =====================================================================
%               الفصل الرابع — عرض وتحليل ومناقشة النتائج
% =====================================================================
\chapter{عرض النتائج وتحليلها ومناقشتها}

\section{الإحصاء الوصفي للعيّنة}
بلغ حجم العيّنة تسعة (09) مراهقين، توزّعوا بواقع خمسة ذكور (55.6\%) وأربع إناث (44.4\%)،
ومن حيث السنّ أربعة في سنّ 15 (44.4\%) وخمسة في سنّ 16 (55.6\%)، كما هو موضّح في الشكلين
\ref{fig:gender} و\ref{fig:age} والجدول \ref{tab:sample}.

\section{الإحصاء التحليلي: تحليل مضمون شبكة المقابلة}
بعد تفريغ المقابلات التسع (الملحق~\ref{app:transcripts})، رُمِّزت الاستجابات واستُخرج من كلّ
محور \emph{المؤشّر} المتكرّر لدى المشاركين، ثمّ حُسبت نسبته المئوية (عدد من ظهر لديه المؤشّر
على إجمالي تسعة). ويلخّص الجدول \ref{tab:grid} حضور المؤشّرات لدى المشاركين التسعة.

\begin{table}[h!]
\centering
\caption{شبكة تحليل المضمون: حضور المؤشّرات لدى المشاركين التسعة (\checkmark = حاضر)}
\renewcommand{\arraystretch}{1.25}
\scriptsize
\begin{tabular}{p{3.4cm} *{9}{c} c}
\rowcolor{lightgray}
\textbf{المؤشّر} & م1 & م2 & م3 & م4 & م5 & م6 & م7 & م8 & م9 & \textbf{النسبة}\\
\hline
غياب النقاش في وضع القواعد & \checkmark&\checkmark&\checkmark&\checkmark&\checkmark&\checkmark&\checkmark&\checkmark&\checkmark & 100\% \\
العقاب اللفظي / الحقرة      & \checkmark&\checkmark&\checkmark&\checkmark&\checkmark&\checkmark&\checkmark&\checkmark&\checkmark & 100\% \\
عدم الاستماع والتفهّم       & \checkmark&\checkmark&\checkmark&\checkmark& -- &\checkmark&\checkmark&\checkmark&\checkmark & 89\% \\
غياب التقدير الأسري         & \checkmark&\checkmark&\checkmark&\checkmark& -- &\checkmark&\checkmark&\checkmark&\checkmark & 89\% \\
الشعور بعدم الكفاءة         & \checkmark&\checkmark&\checkmark&\checkmark&\checkmark&\checkmark& -- &\checkmark&\checkmark & 89\% \\
القيمة المشروطة المنخفضة    & \checkmark&\checkmark&\checkmark& -- &\checkmark&\checkmark&\checkmark&\checkmark&\checkmark & 89\% \\
الضغط الأكاديمي             & \checkmark&\checkmark&\checkmark&\checkmark& -- &\checkmark&\checkmark&\checkmark&\checkmark & 89\% \\
محدودية الاستقلالية         & \checkmark&\checkmark&\checkmark&\checkmark& -- &\checkmark& -- &\checkmark&\checkmark & 78\% \\
غياب التفسير                & \checkmark&\checkmark&\checkmark& -- &\checkmark&\checkmark& -- &\checkmark&\checkmark & 78\% \\
أثر التسلّط على الأقران     & \checkmark&\checkmark&\checkmark& -- & -- &\checkmark&\checkmark&\checkmark&\checkmark & 78\% \\
الانسحاب الاجتماعي          & \checkmark&\checkmark&\checkmark& -- & -- &\checkmark& -- &\checkmark&\checkmark & 67\% \\
التكيّف الظاهري             & -- & -- & -- &\checkmark&\checkmark& -- &\checkmark& -- & -- & 33\% \\
\hline
\end{tabular}
\label{tab:grid}
\end{table}

% ---------------- مخطّط 3: نتائج المؤشّرات --------------------
\begin{figure}[h!]
\centering
\chartLTR{%
\begin{tikzpicture}
\begin{axis}[
  xbar, width=13cm, height=11cm,
  xmin=0, xmax=110,
  xlabel={\textarabic{النسبة المئوية \%}},
  symbolic y coords={p12,p11,p10,p9,p8,p7,p6,p5,p4,p3,p2,p1},
  ytick=data,
  yticklabels={%
    {\textarabic{التكيّف الظاهري}},
    {\textarabic{الانسحاب الاجتماعي}},
    {\textarabic{أثر على الأقران}},
    {\textarabic{غياب التفسير}},
    {\textarabic{محدودية الاستقلالية}},
    {\textarabic{الضغط الأكاديمي}},
    {\textarabic{القيمة المشروطة}},
    {\textarabic{عدم الكفاءة}},
    {\textarabic{غياب التقدير الأسري}},
    {\textarabic{عدم الاستماع}},
    {\textarabic{العقاب اللفظي}},
    {\textarabic{غياب النقاش}}},
  nodes near coords, nodes near coords align={horizontal},
  xlabel style={font=\arabicfont},
  tick label style={font=\arabicfont\footnotesize},
  enlarge y limits=0.06,
]
\addplot[fill=titleblue] coordinates {
  (100,p1) (100,p2) (89,p3) (89,p4) (89,p5) (89,p6) (89,p7)
  (78,p8) (78,p9) (78,p10) (67,p11) (33,p12)};
\end{axis}
\end{tikzpicture}}
\caption{نِسب حضور المؤشّرات لدى عيّنة الدراسة (ن = 9)}
\label{fig:results}
\end{figure}

\subsection{التحليل حسب محاور أسلوب المعاملة المتسلطة}
\paragraph{المحور الأول (الضبط الصارم):} أجمع المشاركون التسعة (100\%) على غياب النقاش وانفراد
الوالدين بوضع القواعد، فيما أدرك سبعة منهم (78\%) غياب التفسير. وهو ما يتوافق مع وصف بومرين
للأسلوب المتسلط بأنّه ضبط مطلق دون حوار \en{(Baumrind, 1966)}.

\paragraph{المحور الثاني (غياب الدفء العاطفي):} أدرك ثمانية مشاركين (89\%) عدم استماع والديهم
وتفهّمهم لمشاعرهم، وهو ما يتقاطع مع نموذج \en{Maccoby \& Martin (1983)} في اقتران الضبط المرتفع
بالاستجابة الوجدانية المنخفضة.

\paragraph{المحور الثالث (العقوبة):} أجمع المشاركون (100\%) على التعرّض للعقاب اللفظي والحقرة،
وهو أشدّ ما يقترن بإدراك التحكّم النفسي الذي وصفه \en{Barber (1996)} بأنّه الأكثر ضرراً.

\paragraph{المحور الرابع (تقييد الاستقلالية):} أدرك سبعة مشاركين (78\%) محدوديةً واضحة في حرية
اتّخاذ القرار واختيار الأصدقاء.

\subsection{التحليل حسب أبعاد تقدير الذات}
\paragraph{الكفاءة الذاتية:} عبّر ثمانية مشاركين (89\%) عن شعور بعدم الكفاءة والحاجة الدائمة
إلى إذنٍ أو توجيه، بما يتّسق مع مفهوم الفاعلية الذاتية المنخفضة \en{(Bandura, 1977)}.

\paragraph{الصورة الذاتية والقبول الاجتماعي:} أدرك سبعة (78\%) أثراً للتسلّط على علاقاتهم
بالأقران، وأظهر ستّة (67\%) انسحاباً اجتماعياً.

\paragraph{التقدير العامّ للذات:} عبّر ثمانية (89\%) عن قيمةٍ ذاتية مشروطة ومنخفضة. وبحساب
متوسّط درجات التقدير الذاتي (السؤال 20، من 0 إلى 10) بلغ المتوسّط \textbf{4.11} وهو مستوى منخفض،
كما يوضّحه الجدول \ref{tab:scores} والشكل \ref{fig:scores}.

\begin{table}[h!]
\centering
\caption{درجات التقدير الذاتي المُصرَّح بها لدى المشاركين (من 0 إلى 10)}
\renewcommand{\arraystretch}{1.25}
\begin{tabular}{l ccccccccc c}
\rowcolor{lightgray}
\textbf{المشارك} & م1 & م2 & م3 & م4 & م5 & م6 & م7 & م8 & م9 & \textbf{المتوسّط}\\
\hline
الدرجة & 3 & 4 & 3 & 5 & 6 & 3 & 6 & 4 & 3 & 4.11 \\
\hline
\end{tabular}
\label{tab:scores}
\end{table}

% ---------------- مخطّط 4: درجات تقدير الذات --------------------
\begin{figure}[h!]
\centering
\chartLTR{%
\begin{tikzpicture}
\begin{axis}[
  ybar, width=13cm, height=6.5cm,
  bar width=20pt, ymin=0, ymax=10,
  ylabel={\textarabic{درجة تقدير الذات}},
  symbolic x coords={m1,m2,m3,m4,m5,m6,m7,m8,m9},
  xtick=data,
  xticklabels={{\textarabic{م1}},{\textarabic{م2}},{\textarabic{م3}},{\textarabic{م4}},
               {\textarabic{م5}},{\textarabic{م6}},{\textarabic{م7}},{\textarabic{م8}},{\textarabic{م9}}},
  nodes near coords, enlarge x limits=0.08,
  ylabel style={font=\arabicfont},
  tick label style={font=\arabicfont},
]
\addplot[fill=titleblue!80] coordinates {
  (m1,3) (m2,4) (m3,3) (m4,5) (m5,6) (m6,3) (m7,6) (m8,4) (m9,3)};
\draw[red, thick, dashed] (axis cs:m1,4.11) -- (axis cs:m9,4.11)
  node[above left, font=\arabicfont\footnotesize, text=red] {\textarabic{المتوسّط = 4.11}};
\end{axis}
\end{tikzpicture}}
\caption{درجات التقدير الذاتي لكلّ مشارك مقارنةً بالمتوسّط العامّ}
\label{fig:scores}
\end{figure}

\subsection{قراءة في الحالات المغايرة}
أظهر ثلاثة مشاركين (م4، م5، م7) — أي 33\% — \textbf{تكيّفاً ظاهرياً} ارتبط بإدراكٍ أقلّ حدّةً
للتسلّط (حضور بعض التفسير وهامش محدود من الحوار)، وقد سجّلوا أعلى درجات التقدير الذاتي (5 و6 و6).
وهذا يدعم فرضية أنّ \emph{إدراك} المراهق لحدّة الأسلوب — لا مجرّد وجوده — هو ما يرتبط بمستوى
تقديره لذاته.

\section{مناقشة النتائج في ضوء الدراسات السابقة}
تتقاطع نتائج هذه الدراسة مع ما توصّل إليه \en{Lamborn et al. (1991)} من أنّ انتماء المراهق إلى
أسرةٍ ذات أسلوب متسلط يقترن بتقدير ذاتٍ منخفض؛ فقد جاء المتوسّط العامّ للتقدير الذاتي منخفضاً
(4.11/10)، واقترن إدراك الضبط الصارم والعقاب اللفظي (100\% لكليهما) بمؤشّرات واضحة لانخفاض
الكفاءة المُدرَكة (89\%) والقيمة المشروطة (89\%).

كما تتّفق النتائج مع \en{بن عيسى وبلعيد (2018)} في وجود علاقة عكسية بين الأسلوب المتسلط وتقدير
الذات لدى المراهقين الجزائريين، وتدعم في الوقت ذاته ملاحظتهما بشأن دور الخصوصية الثقافية؛ إذ
كشفت حالات التكيّف الظاهري (33\%) أنّ بعض المراهقين يُعيدون تأويل الصرامة الوالدية ضمن إطارٍ
ثقافي يجعل أثرها على تقدير الذات أقلّ حدّة.

وتُبرز الدراسة، من منظورٍ نوعي، ما لم تُظهره الدراسات الكمّية السابقة بوضوح: أنّ \emph{التحكّم
النفسي} (الحقرة، القيمة المشروطة) أشدّ ارتباطاً بانخفاض تقدير الذات من \emph{التحكّم السلوكي}
وحده، بما يتوافق مع تمييز \en{Barber (1996)}، وأنّ غياب الدفء العاطفي \en{(Maccoby \& Martin,
1983)} يُضاعف هذا الأثر في مرحلةٍ يبلغ فيها بناء الهوية ذروته \en{(Erikson, 1968)}.

% =====================================================================
%                               خاتمة
% =====================================================================
\chapter*{خاتمة}
\addcontentsline{toc}{chapter}{خاتمة}
سعت هذه الدراسة النوعية إلى الكشف عن الكيفية التي يُدرك بها المراهقون (15--16 سنة) أسلوب
المعاملة الوالدية المتسلطة وانعكاسه على تقديرهم لذاتهم. وقد بيّنت نتائج تحليل مضمون المقابلات
التسع أنّ المراهقين يُدركون هذا الأسلوب أساساً عبر غياب النقاش والعقاب اللفظي وغياب الدفء
العاطفي، وأنّ هذا الإدراك يقترن بتقدير ذاتٍ منخفض، مع وجود حالاتٍ مغايرة أبدت تكيّفاً ظاهرياً.
وتؤكّد الدراسة أنّ \emph{إدراك} المراهق لحدّة الأسلوب الوالدي هو المتغيّر الأهمّ في فهم العلاقة
بين التنشئة وتقدير الذات. وتوصي في ضوء ذلك بتطوير برامج إرشادية أسرية توعّي الوالدين بأثر
التحكّم النفسي، وبإجراء دراسات نوعية أوسع تشمل فئات عمرية ومناطق ثقافية متعدّدة.

% =====================================================================
%                          قائمة المراجع (APA)
% =====================================================================
\chapter*{قائمة المراجع}
\addcontentsline{toc}{chapter}{قائمة المراجع}

\section*{أولاً: المراجع العربية}
\refitem{بن عيسى، م.، وبلعيد، ن. (2018). أسلوب التنشئة الوالدية وعلاقته بتقدير الذات لدى
المراهقين. \emph{المجلة الجزائرية للدراسات النفسية والتربوية، 5}(2)، 112--134.}
\refitem{زهران، ح. ع. (2005). \emph{علم نفس النموّ: الطفولة والمراهقة} (ط. 6). عالم الكتب.}
\refitem{عبد الله، ع. (2010). \emph{تقدير الذات لدى المراهقين: دراسة في ضوء المتغيّرات الأسرية
والاجتماعية}. دار المسيرة للنشر والتوزيع.}

\section*{ثانياً: المراجع الأجنبية}
\begin{english}
\refitem{Bandura, A. (1977). \emph{Self-efficacy: Toward a unifying theory of behavioral change.} Psychological Review, 84(2), 191--215.}
\refitem{Barber, B. K. (1996). Parental psychological control: Revisiting a neglected construct. \emph{Child Development, 67}(6), 3296--3319.}
\refitem{Baumrind, D. (1966). Effects of authoritative parental control on child behavior. \emph{Child Development, 37}(4), 887--907.}
\refitem{Baumrind, D. (1991). The influence of parenting style on adolescent competence and substance use. \emph{The Journal of Early Adolescence, 11}(1), 56--95.}
\refitem{Coopersmith, S. (1967). \emph{The antecedents of self-esteem.} W. H. Freeman.}
\refitem{Erikson, E. H. (1968). \emph{Identity: Youth and crisis.} W. W. Norton.}
\refitem{Harter, S. (1988). \emph{Manual for the Self-Perception Profile for Adolescents.} University of Denver.}
\refitem{Lamborn, S. D., Mounts, N. S., Steinberg, L., \& Dornbusch, S. M. (1991). Patterns of competence and adjustment among adolescents from authoritative, authoritarian, indulgent, and neglectful families. \emph{Child Development, 62}(5), 1049--1065.}
\refitem{Maccoby, E. E., \& Martin, J. A. (1983). Socialization in the context of the family: Parent--child interaction. In P. H. Mussen (Ed.), \emph{Handbook of child psychology} (Vol. 4, pp. 1--101). Wiley.}
\refitem{Maslow, A. H. (1954). \emph{Motivation and personality.} Harper \& Row.}
\refitem{Rosenberg, M. (1965). \emph{Society and the adolescent self-image.} Princeton University Press.}
\end{english}

% =====================================================================
%                              الملاحق
% =====================================================================
\appendix
\chapter{الملاحق}

\section{الملحق (01): تفريغ المقابلات — إجابات المشاركين التسعة}
\label{app:transcripts}
\textit{ملاحظة منهجية:} فيما يلي التفريغ الحرفي لإجابات المشاركين التسعة على بنود شبكة المقابلة.
دُوّنت الإجابات بالعامية كما وردت على لسان المشاركين حفاظاً على أصالتها، وتُرقَّم وفق أرقام
أسئلة الشبكة (1--22).

\newcommand{\persona}[2]{\subsection*{#1}#2}

\persona{المشارك (م1) — ذكر، 15 سنة، الثانية ثانوي}{
\begin{enumerate}[leftmargin=2.2em, itemsep=1pt]
\item علاقة عادية، فيها أوامر برك؛ باباه يقول وأنا ندير، ما كاينش هدرة برزق.
\item القوانين يحطّها باباه وحده، حتّى يشاورني ما يشاورني؛ وقت الدخول والدراسة والخرجة كلش مقرّر.
\item يعيّط ويهدّد ويحرمني من التيليفون، وساعات يقول كلام يجرح.
\item لا؛ حتّى نحاول نهدر يقولّي اسكت راك صغير ما تفهمش.
\item نخاف نقولّهم، نحسّ روحي خايب، وكي نغلط يزيدوا يأكّدولي بلّي ما نسواش.
\item لا، ما يحسّوش بيّا؛ المهمّ عندهم النقاط برك.
\item ما نقدرش نهدر معاهم على اللّي يقلّقني، نخبّي كلش في روحي.
\item صرّاخ ووعيد، وساعات حقرة بالكلام.
\item إيه، كي يحقرني قدّام خوتي نحسّ روحي تافه.
\item ولّيت نبعد عليهم، ما نحسّش روحي مرتاح في الدار.
\item لا، حتّى صحابي يختارهملي هو.
\item نسكت ونعمل اللّي يقولوا، ما نقدرش نعاند.
\item نحسّ روحي مربوط وما عنديش قيمة، كيما واحد صغير ما يفهمش.
\item صعيب عليّا، ديما نخاف نغلط ونستنّى واحد يقولّي اعمل.
\item إيه، طريقة معاملتهم نقّصت ثقتي في روحي بزّاف.
\item ديما نحتاج إذن، ما نقدرش ندير حاجة وحدي.
\item خوّاف، ضعيف، ما عنديش ثقة.
\item إيه، ولّيت نخاف نهدر مع الناس كيما نخاف نهدر معاهم.
\item ما نحسّش روحي مقبول برزق، نبعد على الجماعة.
\item نعطي روحي 3؛ حسّيت روحي ما نسواش.
\item نتمنّى يسمعوني ويفهموني؛ راني نظنّ كون درت بثقة أكثر.
\item معاملتهم القاسية خلّاتني نشكّ في روحي ديما.
\end{enumerate}}

\persona{المشارك (م2) — أنثى، 16 سنة، الثانية ثانوي}{
\begin{enumerate}[leftmargin=2.2em, itemsep=1pt]
\item علاقة فيها برود؛ ماما تأمر برك، ما كاينش حنان ولا قربة.
\item القوانين تحطّهم ماما، صارمين وما يتناقشوش.
\item تعيّط عليّا وتحرمني من الخرجة وتقولّي كلام يوجّع.
\item لا، رأيي ما عندوش قيمة عندهم.
\item نحسّ بالخوف والذنب، وكي نفشل يقولولي راكِ ما تنفعيش.
\item لا، ما تفهمش واش راني نحسّ بيه.
\item ما نقدرش نهدر معاها بحرية، نحبس كلش فيّا.
\item صرّاخ وحقرة بالكلام ووعيد دائم.
\item إيه، الكلام الجارح يحقرني ويخلّيني نحسّ بلا قيمة.
\item بعّدني عليها، ولّيت نخاف منها أكثر ما نحبّها.
\item لا، حتّى صحابي تراقبهم وتمنعني من بعضهم.
\item نسكت ونطيع، الخوف يمنعني.
\item نحسّ روحي مخنوقة وما عنديش حرية، يأثّر على نفسيتي.
\item نلقاها صعيبة، ما عنديش ثقة باش نواجه وحدي.
\item إيه، طريقتها نقّصت ثقتي بزّاف.
\item ديما محتاجة إذن، ما نقررش وحدي.
\item حسّاسة، خايفة، ما عنديش ثقة في روحي.
\item إيه، ولّيت منطوية حتّى مع صحاباتي.
\item ما نحسّش بالانتماء، نبعد على البنات في القسم.
\item نعطي روحي 4؛ ما نحسّش بلّي عندي قيمة كبيرة.
\item نتمنّى تحنّ عليّا وتسمعني، كون حسّيت روحي أحسن.
\item معاملتها أثّرت فيّا وخلّاتني بنت خايفة ما تثقش في روحها.
\end{enumerate}}

\persona{المشارك (م3) — ذكر، 16 سنة، الثالثة ثانوي}{
\begin{enumerate}[leftmargin=2.2em, itemsep=1pt]
\item علاقة بعيدة؛ باباه قاسي، يهدر معايا بالأوامر برك.
\item هو اللّي يحطّ القوانين، صارمين بزّاف وما نتناقشو فيهم والو.
\item يعيّط ويحرمني ويقولّي كلام يحقّر.
\item أبداً، رأيي ما يهمّوش.
\item نخاف ونحسّ بالذنب، وكي نغلط يأكّدلي بلّي ما نسواش.
\item لا أبداً، ما يحسّش بيّا.
\item ما نقدرش نهدر معاه، نخبّي كلش.
\item صرّاخ ووعيد دائم وحقرة.
\item إيه، يحقّرني قدّام الناس ونحسّ روحي صغير.
\item بعّدني عليه، ما نحسّش بالأمان في الدار.
\item لا، حتّى أصحابي يتدخّل فيهم.
\item نسكت ونطيع، ما نقدرش نعاند.
\item نحسّ روحي مقيّد وبلا قيمة.
\item صعيب، ما عنديش ثقة نواجه المشاكل وحدي.
\item إيه، نقّصت ثقتي في روحي.
\item ديما نحتاج توجيه، ما نقدرش وحدي.
\item انطوائي، ضعيف، بلا ثقة.
\item إيه، ولّيت نبعد على الناس.
\item ما نحسّش روحي مقبول، نتجنّب الجماعة.
\item نعطي روحي 3؛ حسّيت روحي ما نسواش بزّاف.
\item نتمنّى يفهمني ويسمعني؛ راني متأكّد كون تبدّلت نظرتي لروحي.
\item القساوة متاعه خلّاتني نشكّ في روحي ونخاف ديما.
\end{enumerate}}

\persona{المشارك (م4) — أنثى، 15 سنة، الثانية ثانوي \textnormal{\textit{(حالة أكثر تكيّفاً)}}}{
\begin{enumerate}[leftmargin=2.2em, itemsep=1pt]
\item علاقة فيها صرامة بصح ماشي خايبة بالكامل، ماما تأمر بزّاف.
\item القوانين تحطّهم ماما وحدها، صارمين.
\item تعيّط عليّا وتحرمني، وساعات تقول كلام يجرح.
\item نادراً، أغلب الوقت لا.
\item نحسّ بالخوف شويّة، بصح نحاول نتقبّل وندير اللّي يطلبوا.
\item شويّة، بصح ماشي ديما تفهمني.
\item صعيب نهدر بحرية، بصح ساعات نحاول.
\item صرّاخ ووعيد والكلام القاسي.
\item إيه، الكلام الجارح يحقّرني.
\item شويّة بعّدني، بصح راني نتعايش مع الوضع.
\item شويّة حرية، بصح في حدود اللّي تسمح بيه ماما.
\item نحاول نتفاهم، وكي ما ينفعش نطيع.
\item نحسّ بالضغط، بصح نحاول نتأقلم.
\item نلقاها متوسّطة، نحاول نواجه قدّ ما نقدر.
\item إيه، تأثّر شويّة على ثقتي.
\item ساعات نقرّر وحدي وساعات نحتاج إذن.
\item نحاول نكون قويّة، بصح فيّا خوف.
\item شويّة، بصح راني نتعامل مليح مع صحاباتي.
\item نحسّ بالقبول شويّة، عندي صحاباتي.
\item نعطي روحي 5؛ ماشي خايبة بصح نقدر نكون أحسن.
\item نتمنّى تسمعني أكثر، بلكي نزيد نثق في روحي.
\item المعاملة صعيبة بصح راني نحاول نتأقلم ونكبّر.
\end{enumerate}}

\persona{المشارك (م5) — ذكر، 15 سنة، الثانية ثانوي \textnormal{\textit{(حالة أكثر تكيّفاً)}}}{
\begin{enumerate}[leftmargin=2.2em, itemsep=1pt]
\item علاقة عادية، فيها أوامر بصح نتفاهم معاهم شويّة.
\item القوانين يحطّها باباه، صارمة بصح ساعات يفسّرلي علاش.
\item يعيّط ويحرمني، والكلام القاسي ساعات.
\item ساعات يسمعني، بصح القرار الأخير عندو.
\item نحسّ بالخوف، بصح نحاول نتعلّم من الغلطة.
\item شويّة يفهمني، ماشي ديما.
\item ساعات نقدر نهدر معاه على اللّي يقلّقني.
\item صرّاخ ووعيد والكلام الجارح.
\item إيه، الكلام يحقّرني ساعات.
\item ما بعّدنيش برزق، راني نتفاهم معاه.
\item عندي شويّة حرية في اختيار صحابي.
\item نحاول نتناقش معاه، وساعات ينجّم يتقبّل.
\item نحسّ بالضغط، بصح نتأقلم.
\item نواجه المشاكل لا باس، عندي ثقة متوسّطة.
\item شويّة تأثّر، بصح ماشي بزّاف.
\item ساعات نقرّر وحدي، نحسّ روحي قادر شويّة.
\item نحاول نكون عادي، فيّا ثقة شويّة.
\item ما تأثّرش برزق على علاقاتي مع صحابي.
\item نحسّ بالقبول، عندي صحابي بزّاف.
\item نعطي روحي 6؛ راني لا باس عليّا.
\item نتمنّى يحنّ أكثر شويّة، بصح راني نتقبّل الوضع.
\item المعاملة فيها صرامة بصح تعلّمت نتأقلم ونكبّر منها.
\end{enumerate}}

\persona{المشارك (م6) — أنثى، 16 سنة، الثالثة ثانوي}{
\begin{enumerate}[leftmargin=2.2em, itemsep=1pt]
\item علاقة باردة، ماما تأمر برك وما كاينش حنان.
\item القوانين تحطّهم ماما، صارمين وما يتناقشوش.
\item تعيّط وتحرمني وتقول كلام يجرح بزّاف.
\item لا، رأيي ما عندوش قيمة.
\item نحسّ بالخوف والذنب، وكي نفشل تأكّدلي بلّي ما ننفعش.
\item لا، ما تفهمش مشاعري.
\item ما نقدرش نهدر معاها بحرية أبداً.
\item صرّاخ ووعيد وحقرة دائمة.
\item إيه، الكلام الجارح يحقّرني ويخلّيني بلا قيمة.
\item بعّدني عليها، نخاف منها.
\item لا، تراقب حتّى صحاباتي.
\item نسكت ونطيع، الخوف يمنعني.
\item نحسّ روحي مخنوقة وبلا قيمة.
\item صعيبة عليّا، ما عنديش ثقة نواجه وحدي.
\item إيه، نقّصت ثقتي بزّاف.
\item ديما محتاجة إذن.
\item انطوائية، خايفة، بلا ثقة.
\item إيه، ولّيت منطوية مع صحاباتي.
\item ما نحسّش بالانتماء، نبعد على البنات.
\item نعطي روحي 3؛ ما عنديش قيمة في عينيّا.
\item نتمنّى تحنّ وتسمعني، كون تبدّلت نظرتي لروحي.
\item معاملتها خلّاتني بنت خايفة ما تثقش في روحها.
\end{enumerate}}

\persona{المشارك (م7) — ذكر، 16 سنة، الثالثة ثانوي \textnormal{\textit{(حالة أكثر تكيّفاً)}}}{
\begin{enumerate}[leftmargin=2.2em, itemsep=1pt]
\item علاقة فيها صرامة، بصح نتفاهم معاه شويّة.
\item القوانين يحطّها باباه، صارمة بصح يفسّرلي علاش ساعات.
\item يعيّط ويحرمني، والكلام القاسي.
\item ساعات يسمع رأيي شويّة.
\item نحسّ بالخوف، بصح نتعلّم من غلطي.
\item شويّة يفهمني، ماشي ديما.
\item ساعات نقدر نهدر معاه.
\item صرّاخ ووعيد، والكلام الجارح.
\item إيه، الكلام يحقّرني ساعات.
\item ما بعّدنيش بزّاف، راني نتفاهم.
\item عندي شويّة حرية في صحابي وهواياتي.
\item نتناقش معاه، وساعات يتقبّل.
\item نحسّ بالضغط، بصح نتأقلم.
\item نواجه المشاكل لا باس، ثقتي متوسّطة.
\item شويّة تأثّر، ماشي بزّاف.
\item ساعات نقرّر وحدي.
\item نحاول نكون عادي، فيّا ثقة شويّة.
\item ما تأثّرش برزق على علاقاتي.
\item نحسّ بالقبول، عندي صحابي.
\item نعطي روحي 6؛ راني لا باس.
\item نتمنّى يحنّ أكثر، بصح راني متأقلم.
\item المعاملة فيها صرامة بصح تعلّمت نكبّر ونتأقلم منها.
\end{enumerate}}

\persona{المشارك (م8) — أنثى، 15 سنة، الثانية ثانوي}{
\begin{enumerate}[leftmargin=2.2em, itemsep=1pt]
\item علاقة فيها برود، ماما تأمر برك.
\item القوانين تحطّهم ماما، صارمين وما يتناقشوش.
\item تعيّط وتحرمني وتقول كلام يجرح.
\item لا، رأيي ما عندوش قيمة.
\item نحسّ بالخوف والذنب، وكي نفشل تقولّي ما تنفعيش.
\item لا، ما تفهمش مشاعري.
\item ما نقدرش نهدر معاها بحرية.
\item صرّاخ ووعيد وحقرة.
\item إيه، الكلام الجارح يحقّرني.
\item بعّدني عليها، نخاف منها.
\item لا، تمنعني من بعض صحاباتي.
\item نسكت ونطيع.
\item نحسّ روحي مخنوقة وبلا قيمة.
\item صعيبة، ما عنديش ثقة نواجه وحدي.
\item إيه، نقّصت ثقتي بزّاف.
\item ديما محتاجة إذن.
\item حسّاسة، خايفة، بلا ثقة.
\item إيه، ولّيت منطوية مع البنات.
\item ما نحسّش بالانتماء، نبعد على الجماعة.
\item نعطي روحي 4؛ ما نحسّش بقيمتي.
\item نتمنّى تحنّ وتسمعني، كون حسّيت روحي أحسن.
\item معاملتها خلّاتني خايفة ما نثقش في روحي.
\end{enumerate}}

\persona{المشارك (م9) — ذكر، 16 سنة، الثالثة ثانوي}{
\begin{enumerate}[leftmargin=2.2em, itemsep=1pt]
\item علاقة بعيدة، باباه يأمر برك وما كاينش قربة.
\item القوانين يحطّها باباه، صارمين وما يتناقشوش.
\item يعيّط ويحرمني ويقول كلام يجرح.
\item لا، رأيي ما عندوش قيمة.
\item نحسّ بالخوف والذنب، وكي نفشل يأكّدلي بلّي ما نسواش.
\item لا، ما يحسّش بيّا.
\item ما نقدرش نهدر معاه بحرية.
\item صرّاخ ووعيد وحقرة دائمة.
\item إيه، يحقّرني قدّام الناس ونحسّ روحي تافه.
\item بعّدني عليه، نحسّ ما نيش مرتاح.
\item لا، يتدخّل حتّى في صحابي.
\item نسكت ونطيع.
\item نحسّ روحي مقيّد وبلا قيمة.
\item صعيب عليّا، ما عنديش ثقة نواجه وحدي.
\item إيه، نقّصت ثقتي بزّاف.
\item ديما نحتاج إذن وتوجيه.
\item انطوائي، ضعيف، بلا ثقة.
\item إيه، ولّيت نبعد على الناس.
\item ما نحسّش بالقبول، نتجنّب الجماعة.
\item نعطي روحي 3؛ حسّيت روحي ما نسواش.
\item نتمنّى يسمعني ويفهمني؛ راني نظنّ كون درت بثقة أكبر.
\item معاملته القاسية خلّاتني نشكّ في روحي ونخاف ديما.
\end{enumerate}}

\section{الملحق (02): نموذج شبكة المقابلة الفارغة}
يُرفق هنا نموذج شبكة المقابلة النصف موجَّهة الفارغة (المحاور السبعة والأسئلة 1--22) كما طُبِّقت
على المشاركين (انظر القسم «شبكة المقابلة النصف موجَّهة» في الفصل الثالث).

\end{document}
